\subsection{Registers and the precise claims}

The nine logical registers, in order, are
\[
 (r_0,r_1,r_2,r_3,r_4,r_5,r_6,r_7,r_8)
       =(N,T,W,A,D,H,S,R,L).
\]
Here $N$ is the fresh-candidate counter, $L$ the continuation code,
$T,W,H$ the displayed theorem, theorem stack, and cached theorem, and $D$
an arithmetic operand offset. $A,S,R$ are work registers. Source site
numbers in this note are zero-based indices in the delivered instruction
array.

The standard clean scheduler entry at source site 25 is
\[
 (N,T,W,0,D,H,0,0,L).
\tag{1}
\]
Its abstract scheduler state is $\langle N,L,\sigma\rangle$, where
$\sigma=\langle T,W,D,H\rangle$. A store is \emph{sound} if $T$ and $H$
have derivations in the source numerical calculus and every live canonical
entry of $W$ has such a derivation. There is no soundness condition on
$D$. In particular $\langle0,0,D,0\rangle$ is sound for every $D$ because
code zero has a fixed derivation.

\begin{theorem}[Arbitrary clean scheduler seed]\label{thm:clean}
For all natural numbers $N,L,D$, the source computation from site 25 with
registers
\[
                (N,0,0,0,D,0,0,0,L)
\tag{2}
\]
halts if and only if the numerical calculus has a finite refutation.
More generally the same equivalence holds with any sound store in (1).
The proof is formalizable in PA uniformly in all the displayed numbers.
In particular neither $N\ge2$ nor $L=0$ is required.
\end{theorem}

\begin{theorem}[Dirty work registers can be normalized]\label{thm:dirty}
Fix any source site $j$ with $25\le j\le44$. From that site let
\[
          (N,T,W,A,D,H,S,R,L)=(N,0,0,A,D,0,S,R,L),
\tag{3}
\]
where all six values $N,A,D,S,R,L$ are arbitrary natural numbers.
The source computation halts if and only if the numerical calculus has a
finite refutation. This equivalence too is provable in PA by the finite
source-computation and syntactic proof arguments below.
\end{theorem}

Under the inherited PA calibration, either theorem yields
\[
 \PA\vdash \bigl(\neg\Halt(\text{source from the stated seed})
                  \leftrightarrow\Con(\PA)\bigr).
\]
Theorems~\ref{thm:clean}--\ref{thm:dirty} do not enlarge the permitted
physical tape-frame format. A caller must still check the marker, encoded
source path, head position, padding, complete register-file encoding, and
spare-register tail required by its literal runtime certificate.

\subsection{Cofinality from every scheduler state}

Write $p(a,b)=(a+b)(a+b+1)/2+b$ and let \texttt{parse} be the command
decoder. The unchanged scheduler uses
\[
 \operatorname{seed}(N,L)=
 \begin{cases}N&L=0,\\L-1&L>0,\end{cases}
 \qquad
 \operatorname{nextN}(N,L)=
 \begin{cases}N+1&L=0,\\N&L>0.\end{cases}
\tag{4}
\]
The next continuation is
$\operatorname{parse}(\operatorname{seed}(N,L)).\mathrm{next}$.
The component bounds for Cantor decoding give
\[
                 \operatorname{parse}(u).\mathrm{next}\le u.
\tag{5}
\]
Consequently a nonzero continuation decreases strictly at each
nonhalting command. From an arbitrary $\langle N,L,\sigma\rangle$,
either a halt occurs or at most $L$ commands reach a state with continuation
zero and the same fresh counter $N$.

A fresh step from $\langle n,0,\sigma\rangle$ increments the counter to
$n+1$ and leaves a continuation at most $n$. Thus at most $n+1$ additional
commands either halt or reach fresh counter $n+1$ with continuation zero.
For each $K\ge N$, a state with fresh counter $K$ and continuation zero
is reached unless a halt occurs first, within
\[
        L+\sum_{n=N}^{K-1}(n+1)
       =L+\frac{(K-N)(K+N+1)}2
\tag{6}
\]
commands. The theorem does not assume that the unsuccessful trials reset
the proof store.

\subsubsection{Padding a halting command word above the initial counter}

For a command $c$ define
\[
\begin{aligned}
 \operatorname{ec}(\mathsf{Shift},t)&=p(0,t),\\
 \operatorname{ec}(\mathsf{MP}_x,t)&=p(1,p(x,t)),\\
 \operatorname{ec}(\mathsf{AX}_{a,x,y},t)
      &=p(a+2,p(y,p(x,t))).
\end{aligned}
\]
Define a finite command-word encoding by
\[
 \operatorname{ew}([],t)=t,\qquad
 \operatorname{ew}(c::w,t)=\operatorname{ec}(c,\operatorname{ew}(w,t))+1.
\tag{7}
\]
The Cantor component bounds prove $t\le\operatorname{ew}(w,t)$.
Parsing the positive continuation in (7) executes the indicated command
and then retains precisely the encoded remaining word.

Let $q$ be a numerical-calculus refutation. The inherited command compiler
provides a nonempty word $w_q$ which halts from \emph{every} incoming store.
For an initial fresh counter $N$, choose
\[
       t=N+1,\qquad C=\operatorname{ew}(w_q,t),\qquad K=C-1.
\tag{8}
\]
Then $C\ge N+1$ and $K\ge N$. By (6), the scheduler either has already
halted or reaches $\langle K,0,\tau\rangle$ for some store $\tau$.
Its next transition is identical to that from
$\langle C,C,\tau\rangle$: both parse $C-1$ and both retain fresh counter
$C$. The latter state executes $w_q$ before reading its tail, and hence
halts. This proves completeness from every $N,L,\sigma$, even without a
soundness hypothesis on the initial store.

For soundness, retain the existing proof annotation for $T,H,W$.
Every successful command transforms those finite annotations into
annotations for the output store. In AX the inequality needed by the
restricted propositional head is $D\le D+x$, valid for every natural $D$.
An MP success obtains its new theorem from the two annotated premises.
At a halt, the tested proof slot yields a refutation. Thus soundness needs
only the stated proof-store invariant, not a bound on $N,L,$ or $D$.
This proves Theorem~\ref{thm:clean}.

\subsection{The literal normalization prefix}

The HAZ source instructions 25 through 47 agree exactly with the same
slice of the primary PA194 instruction array. In these instructions
\texttt{dec r z p} decrements register $r$ and jumps to $p$ when it is
positive, and otherwise jumps to $z$ without changing it.

\begin{center}\small
\begin{tabular}{r l @{\qquad} r l}
\toprule
Site&Instruction&Site&Instruction\\
\midrule
25&\texttt{dec 8 27 26}&37&\texttt{jump 41}\\
26&\texttt{jump 34}&38&\texttt{dec 7 41 39}\\
27&\texttt{dec 0 31 28}&39&\texttt{inc 6; goto 40}\\
28&\texttt{inc 8; goto 29}&40&\texttt{jump 38}\\
29&\texttt{inc 6; goto 30}&41&\texttt{jump 34}\\
30&\texttt{jump 27}&42&\texttt{dec 7 45 43}\\
31&\texttt{inc 0; goto 32}&43&\texttt{inc 8; goto 44}\\
32&\texttt{dec 6 34 33}&44&\texttt{jump 42}\\
33&\texttt{jump 31}&45&\texttt{dec 6 48 46}\\
34&\texttt{dec 8 42 35}&46&\texttt{inc 3; goto 47}\\
35&\texttt{inc 7; goto 36}&47&\texttt{jump 45}\\
36&\texttt{dec 6 38 37}&&\\
\bottomrule
\end{tabular}
\end{center}

\begin{lemma}[Prefix termination from arbitrary work-register values]
Starting at any site 25 through 44, the displayed block reaches site 48
in finitely many steps. Registers $T,W,D,H$ remain unchanged. At site 48
both $S$ and $R$ are zero; the values of $N,A,L$ are natural numbers, with
no additional restriction needed below.
\end{lemma}
\begin{proof}
At sites 27--30, each complete loop decreases $N$ and increases $L,S$.
It reaches site 31 when $N=0$. At sites 31--33, each complete loop
decreases $S$ and increases $N$, until site 34 is reached with $S=0$.
Starting in the middle of either loop adds a finite partial iteration
and does not affect this termination argument. Site 25 either takes the
two-step continuation branch to 34 or enters the preceding drain/restore
loops; site 26 jumps directly to 34. No hypothesis that $A,S,R$ initially
vanish was used.

At site 34, every complete passage through sites 35--41 decreases $L$
by one. The only inner loop, sites 38--40, decreases $R$ and increases
$S$, so it terminates from every value of $R$. Therefore induction on
the value of $L$ proves that sites 34--41 reach site 42. Again any entry
in the middle of this block first completes a finite partial passage.

Sites 42--44 drain $R$ into $L$, reaching site 45 with $R=0$.
Sites 45--47 then drain $S$ into $A$, reaching 48 with $S=R=0$.
An initial entry at 42, 43, or 44 merely changes the first finite partial
passage of the first drain. Throughout, no instruction writes $T,W,D,H$.
\end{proof}

This argument deliberately does not extend the arbitrary-$R$ assertion
to initial sites 45--47, because those sites do not clear an incoming
nonzero $R$. Nor does it assert that arbitrary values of the proof-store
registers are sound.

\subsubsection{Why the normalized selector is a legitimate command}

At site 48, write the resulting register values as
\[
               (N',0,0,a',D,0,0,0,\ell').
\tag{9}
\]
There is no need for $a',\ell'$ to be the intended decoding of the
original startup's $L$. They are the exact Cantor decoding of
$u=p(a',\ell')$. Consequently (9) is precisely the selector state
produced by an ordinary clean fetch of $u$, with the same proof store
and fresh counter $N'$. Equivalently, a clean scheduler state
$\langle N',u+1,\langle0,0,D,0\rangle\rangle$ reaches (9) after its
positive-continuation fetch.

Both the dirty-entry normalization and that clean positive-continuation
fetch have finite, nonhalting prefixes ending in the identical complete
source configuration (9). Determinism therefore identifies their future
halting behavior. Apply Theorem~\ref{thm:clean} to the hypothetical clean
scheduler entry with counter $N'$, continuation $u+1$, and sound store
$\langle0,0,D,0\rangle$. This proves Theorem~\ref{thm:dirty} directly,
without an additional assumption about the first selected command.

\subsection{Internalization and physical register slots}

The proofs above use induction on natural counters, finite command words,
finite derivations, and literal computation traces. Each loop rank is an
explicit decreasing natural number. The cofinality bound (6) is a
primitive-recursive bound on scheduler transitions. Encoding (7), the
choice (8), and all proof annotations are syntactic primitive-recursive
constructions. Thus PA verifies the arguments uniformly; they are not
restricted to externally standard initial counters. No truth predicate
for PA and no soundness or reflection principle for PA is invoked.
This note supplies an ordinary mathematical PA argument, not a newly
compiled end-to-end Lean theorem.

For the physical allocation used in the HAZ packet, the map from logical
register indices to physical register slots is
\[
                    [6,2,0,3,4,7,8,5,1].
\]
The primary PA194 allocation has instead
$[6,2,0,4,3,7,8,5,1]$; HAZ exchanges the $A,D$ slots. Both registers
are unrestricted in the broad seed criterion, so the required zero
slots below are the same for both allocations.
Accordingly the broad sufficient condition (3) imposes exactly
\[
 \text{physical slot }2=0,\qquad
 \text{physical slot }0=0,\qquad
 \text{physical slot }7=0.
\]
The other six working slots may be arbitrary. This statement assumes a
valid complete runtime frame for one of the stated source sites. In
particular, native control 84 with the path of site 25 is not by itself
enough without the tape-format conditions.

\subsection{Optional broader stores certified directly in PA}

For the requested PA consistency equivalence, the initial annotations may
be PA proofs of the translated formulas. They need not be derivations in
the numerical calculus itself. Completeness still uses the all-store
halting word, and soundness carries those initial PA proofs through the
same syntactic command templates. This yields a further sufficient test
for candidate startups without changing their code.

The accompanying \texttt{seed\_truth.py} is a deliberately partial
syntactic recognizer. Its output is \emph{true}, \emph{false}, or
\emph{unknown}; true and false come with a finite tree of the following
proof templates. Here a true certificate for $B$ means a PA proof of its
translation, and a false certificate means a PA proof of its negation,
in both cases universally over the remaining free variables.
\begin{center}\small
\begin{tabular}{p{0.51\linewidth}p{0.35\linewidth}}
\toprule
Recognized form or premises&Certified value\\
\midrule
Code zero&True, by the truth template\\
$E(a,a)$&True, by equality reflexivity\\
$M(a,a)$&False, by finite-bit irreflexivity\\
$I(B,B)$&True, by implication identity\\
$I(B,C)$ with $B$ false&True, by a false antecedent\\
$I(B,C)$ with $C$ true&True, by a true consequent\\
$I(B,C)$ with $B$ true and $C$ false&False\\
$U(a,B)$ with $B$ true&True, by generalization\\
$U(a,B)$ with $B$ false&False, by specialization at zero\\
\bottomrule
\end{tabular}
\end{center}
The false atomic certificate is justified by the PA theorem
$\operatorname{Bit}(s,z)\Rightarrow z<s$, applied with $z=s$.
All recursion uses strictly smaller decoded formula children. These
rules therefore yield primitive-recursive proof-producing templates;
they are not a full arithmetic truth predicate. An unknown result makes
no claim about provability or truth.

Apply this recognizer to $T,H$ and every canonical decoded entry of
$W$, stopping at numerical stack zero. If all receive true certificates,
the dirty-work-register normalization theorem has the same direct PA
consistency conclusion. For example $T,H$ may lie in $\{0,2,3,4\}$:
these are truth, $\forall x_0\top$, $x_0=x_0$, and
$M(0,0)\to\top$. Sound stack values include $0$, $3=p(2,0)$,
$6=p(3,0)$, and $10=p(4,0)$. By contrast $4=p(1,1)$ contains a
contradiction head and is rejected.

This optional criterion proves a direct PA-refutation equivalence. It
should not be misstated as a source-numerical-calculus equivalence unless
the chosen initial annotations are also shown to be derivations in that
calculus. The script checks only these proof-store certificates and the
source-site interval, not the literal tape-frame format.

\subsection{Source provenance}

Repository snapshot:
\url{https://github.com/VladimirReshetnikov/ProveIt/tree/5784600f8709872fa258af66a94df8deb7075c2f}.
The following files are relative to
\file{Computability/BusyBeaver/Research/PA/}:
\begin{itemize}
\item \file{pa194-command-source/PA194CommandRegisters.lean} gives the
      exact logical register map.
\item \file{pa194-command-source/PA194FetchDispatch.lean} defines
      \texttt{seed} and \texttt{nextN}.
\item \file{pa194-scheduler-fairness/PA194Fairness.lean} proves
      \texttt{parse\_next\_le}, \texttt{drain}, \texttt{reach\_fresh},
      \texttt{encodeWord\_tail\_le}, and \texttt{fresh\_same\_advance}.
      These are already stated for arbitrary states or counters.
\item \file{pa194-scheduler-fairness/PA194EncodedWord.lean} proves
      \texttt{derivation\_halts\_scheduler} for every tail, counter,
      and incoming store.
\item \file{pa194-calculus-soundness/PA194CalculusSoundness.lean}
      defines \texttt{SoundStore} and proves command and scheduler
      soundness for every sound store.
\item \file{pa194-command-source/PA194FetchPrefixBinding.lean},
      \texttt{PA194FetchPrefix.lean}, and \texttt{PA194FetchDecode.lean}
      bind the fetch blocks to literal source sites.
\end{itemize}

The primary instruction slice is
\file{docs/pa194_conpa_2026-11-13/construction/source/pa194.instructions.json}.
Its sites 25--47 were compared exactly with
\file{machine/pa194haz.instructions.json} extracted from
\file{docs/incoming/pa194haz_research_2026-11-14.zip}.
The accompanying JSON receipt records the checked slice and source paths.
For a cached-head variant, its existing all-store command compiler and
soundness proof replace the original command formulas; the fetch and
cofinality proofs are identical.
